\documentclass[]{../../../../tex/classes/styledArticle}
\usepackage{../../../../tex/packages/hebrewSupport}


\author{שחר פרץ}
\title{\textit{היסטוריה 50}}
\begin{document}
	\maketitle
	\section*{הדרוזים והצ'רקסים}
	לממההה למההה מורים עושים מצגות עם NoteBookLM זה ממש רואים שזה AI
	
	הדרוזים והצ'רקסים הם \textit{קבוצות אתניות} – לא קבוצות לאום, משום שהם שאינם שואפים לעצמאות מדינית. 
	
	בהקשר של הדרוזים, מדובר על עדה, שנמצאת ברובה בשטחה שלנו, ונאמנותם נתונה למדינה בהם הם מצויים, אבל גם לעדה – לאחרונה הדרוזים בישראל חצו את הגבול ועזרו לדרוזים בסוריה. 
	
	ב־1948, בהתחלה, הדרוזים (באופן כללי) ''לא היו סגורים על עצמם`` – אין פה מדינת ישראל עדיין. היו קבוצות בתוך החברה הדרוזית שהתלבטו מה לעשות (זכרו, מדובר לפני קום המדינה). 
	
	לגבי הדרוזים – מדובר בדת מונוטאיסטית סגורה. מעט מאוד דרוזים מכירים את הדת שלהם, רק השכבה העליונה מכירה את עיקרי הדת. אם מישהו מהדרוזים מחליט להתחתן מחוץ לעדה, הוא יוצא מהעדה, וילדיו לא נחשבים דרוזים. זו ממש בעיה בתוך העדה הדרוזית. הם מאמינים בגלגול נשמות. עבורם המקום הקדוש הוא קבר הנביא שועייב. 
	
	לגבי הצ'רקסים – מקורם בצפון הקווקז, שבכלל לא קרוב אלינו. דתם מוסלמים סוניים, אך הם מבדילים עצמם משאר הסונים. מצויים בשני כפרים בגליל – כפר כמא וריחאניה. גם עם הדרוזים היו יחסי שכנות טובים, אבל הצ'רקסים לא התלבטו כלל. 
	
	על הדרוזים הופעלו לחצים, גם מהצבאות הלא סדירים הערבים, וגם מגופים אחרים, לחבור אליהם. מנגד היה לחץ מהציונים שיחברו אליהם. 
	
	קרב רמת יוחנן הוא נקודת המפנה כאן. ''צבא ההצלה`` של הקצין שכיב והאב, היה עם יחידה דרוזית, רובם לא דרוזים מישראל (שהיו ניטרלים בעיקרים) אלא מבחוץ (סוריה). בקרב הזה היו מספר קרבות מול קיבוץ רמת יוחנן. האבדות והתבוסה הדרוזית הובילו לנקודת מפנה – חלק חזרו לסוריה, חלק עברו לחיות פה, והעדה הדרוזית שחיה בארץ ישרלא החליטה לקשור את קשרה עם מדינה ישראל ולהיות חלק ממנה. 
	
	ב־1948 הוקמה יחידה ששמה ''יחידת המיעוטים``, שעד לא מזמן גם הייתה קיימת. מי שהיו בה הם דרוזים וצ'רקסים, וקצת מאוחר יותר, הבדואים. היום האוכלוסיות האלו מפוזרות בין כל היחידות בצה''ל, בהתחלה הם היו רק ביחידה הזו. עם זאת הבדואים עד היום מנותבים ליחידת הגששים, בגלל כישוריהם. במבצע חירם בו השתתפו גם הצ'רקסים, שהתקיים במהלך ההפוגה השנייה בהחלק השני של מלחמת העצמאות, והתרחש בגליל העליון ומטרתו הייתה לכבוש אותו, יחידת המיעוטים נלחמה שם. היה ניסיון שיעשו איזשהו תיאום, שלא צלח, ונהרגו שם 14 אנשים (רובם דרוזים). 
	
	בהתחלה הייתה דרישה לפרק את הדרוזים מנשקם (היה להם נשק כמו הרבה אנשים פה בארץ). ב־1949 בן־גוריון ביטל את הכוונה הזו וביקש להכיר בהם קקהילה אוטונומית. הדרוזים והצ'רקסים, בניגוד לבדואים, מחוייבים בגיוס חובה, שהחל ב־1956. אחוזי הגיוס שלהם מאוד גבוהים, בקרב הגברים בלבד. ב־1957 הם הוכרו כעדה דתית נפרדת בישראל (גם השיפוט הדתי שלהם נמצא באוטונומיה מלאה). 
	
	מדינת ישראל לא משקיעה בכפרים הדרוזים בתשתיות, ויש פער מטורף בין ההשקעה של העדה לבין ההשקעה של מדינת ישראל. חוק הלאום עורר של מחאה גדולה. 
	
	לבגרות, חשוב להכיר את מבצע חירם. לפני שנה או שנתיים, הייתה שאלה על הדרוזים בה בקשו קרב שמבטא את שיתוף הפעולה של הדרוזים. ואז כל תלמידי מדינת ישראל כתבו קרב רמת יוחנן. זה לא נכון. זה ההפך הגמור. בקרב הזה הדרוזים תקפו את מדינת ישראל. התשובה הנכונה היא מבצע חירם, במקום הראשון בו רואים את הדרוזים כחלק מהכוח היהודי. הצ'רקסים גם השתתפו במבצע חירם, וגם במבצע דקל. אפשר לדבר על קרב רמת יוחנן כדי לדבר על השינוי והמפנה, אך אם מבקשים קרב שמבטא את הקשר, בקרב רמת יוחנן אפשר לדבר על הקשר רק כנקודת מפנה בקשר. 
	
	''נא לא לכתוב את זה שחר, שהמורה אמרה שהיא פושעת.``
	
	הספרים יחסית ישנים והנושא של הדרוזים נכנס רק אח''כ לתוכנית הלימודים, ולכן זה לא מופיע בספרים. 
	
	עוד הערה לגבי הבחינה. בפרק על הכרזת העצמאות, אין מקור לבעד/נגד ההכרזה במועד יציאת הבריטים, בפרק עצמו. אבל עמוד לפני יש את מפת מלחמת העצמאות במלחמה הפנימית, מדצמבר עד מאי. זה טיעון בעד. אפשר לראות את ההתרחבות של היישוב היהודי. אפשר גם לקחת את זה לירושלים הנצורה ולטעון נגד. קודם לכן יש גם תמונה של חיילים שבויים בעיר העתיקה. בעמוד 92 יש את הפקודה להקמת צה''ל, שהסעיף השני שלה מדבר על גיוס חובה במצב חירום. באופן דומה בפרק על שלילת הדמוקרטיה של יגאל משעול בספר של שואה, אין שום מקור. אבל יש המון מקורות לנושא הזה בפרקים אחרים. 
	
	בכל מקרה, מה שמבקשים, אכן קיים, בכל ספר. כך גם במתכונות של בית הספר. גם בנושא הניסוחים, בחינת הבגרות ניסויים ונבדקת ע''י מכון סולד. בטח ובטח בעניין הספרים – משרד החינוך עובר ומוודא שבכל ספר יש מקור אחד. 
	
\end{document}